\newpage
\selectlanguage{greek}
\chapter*{Περίληψη}
\addcontentsline{toc}{chapter}{Περίληψη}

Αυτή η διπλωματική εργασία παρουσιάζει το \textlatin{\textbf{S\textsuperscript{4}D} (Style and Structure Similarity for Site Domains)}, μια ερευνητική πλατφόρμα για την ανάλυση ομοιοτήτων μεταξύ ιστοσελίδων βάσει της δομής και του οπτικού τους στυλ. Το σύστημα επεξεργάζεται ιστοσελίδες αναλύοντας τη δομή του \textlatin{DOM} και τις ιδιότητες \textlatin{CSS}, συνδυάζοντάς τα σε μία ενιαία μορφή για συγκριτική αξιολόγηση.

Οι χρήστες μπορούν να εισάγουν οποιονδήποτε αριθμό \textlatin{URLs} για ανάλυση. Κάθε σελίδα επεξεργάζεται ανεξάρτητα για την εξαγωγή δομικών και στυλιστικών χαρακτηριστικών. Κατά τη φάση υπολογισμού της ομοιότητας, οι συγκρίσεις πραγματοποιούνται σε επίπεδο \textlatin{domain}, ανεξάρτητα από τον τρόπο ομαδοποίησης των \textlatin{URLs}. Το σύστημα υποστηρίζει την ανάλυση σε διαφορετικά βάθη \textlatin{DOM} (επίπεδα 1 έως 10).

Η ομαδοποίηση των στοιχείων πραγματοποιείται με τον αλγόριθμο \textlatin{HDBSCAN}, ενώ η ποιότητα των ομάδων αξιολογείται με το \textlatin{DBCV}. Τα απομονωμένα στοιχεία επανατοποθετούνται σε κατάλληλες ομάδες όταν είναι δυνατόν. Η πλατφόρμα επιτρέπει τη διεξαγωγή πολλαπλών πειραμάτων με διαφορετικά σύνολα \textlatin{URLs} και παραμέτρους, καλύπτοντας ποικίλες ερευνητικές ανάγκες.

Το \textlatin{\textbf{S\textsuperscript{4}D}} υλοποιείται με \textlatin{React} για το περιβάλλον χρήστη, \textlatin{Flask} για τις διεπαφές \textlatin{API} και \textlatin{Python} για την επεξεργασία δεδομένων. Ενσωματώνει το \textlatin{Gnuplot} για παραγωγή γραφημάτων και αρχείων κατάλληλων για ακαδημαϊκή χρήση. Ένα αρχείο αποτελεσμάτων επιτρέπει τη σύγκριση διαφορετικών πειραμάτων και την παρακολούθηση της προόδου.

Το σύστημα εξάγει \textlatin{HTML} σελίδες με χρωματική κωδικοποίηση που αναδεικνύουν τις δομικές ομοιότητες, καθώς και πίνακες συσχετίσεων βασισμένους σε υπολογισμούς συνημιτόνου. Οι χρονικές μετρήσεις συμβάλλουν στην αξιολόγηση της απόδοσης του συστήματος. Οι δοκιμές δείχνουν ότι το \textlatin{\textbf{S\textsuperscript{4}D}} συγκρίνει και ομαδοποιεί με ακρίβεια στοιχεία ιστοσελίδων σε διαφορετικά \textlatin{domains}, προσφέροντας ένα αξιόπιστο εργαλείο για δομική και οπτική ανάλυση.

Η πλήρης υλοποίηση του συστήματος είναι διαθέσιμη στο: \textlatin{\url{https://github.com/pfavvatas/lib_url_to_img}}
